\subsubsection{The higher-genus three-point example: horizontal wild inertia}

\begin{proposition}[Proved under stated hypotheses: one horizontal tame condition]
\label{D:Schroer-horizontal-tame}
Let $k$ be algebraically closed, let $C/k$ be a smooth proper
connected curve, and choose distinct points $b_0,b_1,b_2\in\mathbb P^1(k)$
and distinct points $c_1,c_2\in C(k)$.  Let
\[
 X=\operatorname{Bl}_{(b_0,c_1),(b_1,c_1),(b_2,c_2)}
       (\mathbb P^1\times C).
\]
Write $R_i$ for the strict transform of $\{b_i\}\times C$, $E_i$
for the corresponding exceptional curve, and $A$ for the strict
transform of $\mathbb P^1\times\{c_1\}$.  Put
\[
 U=X\setminus(A\cup R_0\cup R_1\cup R_2).
\]
Every finite \'etale cover of $U$ which is tame at the divisorial
valuation defined by $A$ is tame at each of the valuations defined
by $R_0,R_1,R_2$.

Consequently, if $P_A$ denotes a wild inertia group at $A$, the
natural quotient map induces an isomorphism
\[
 \pi_1(U)/\overline{\langle\!\langle P_A\rangle\!\rangle}
 \simeq Q(X,X\setminus U).
\]
In particular this quotient is topologically finitely generated.
\end{proposition}

\begin{proof}
We may pass to connected components and then to a finite Galois
closure.  Tameness is preserved under passage to this closure,
composita and subextensions; see
\cite[section 15.116, Lemmas 15.116.5--9]{StacksBIndexTame}.
Thus let $V\to U$ be a connected finite Galois cover with group $G$.
Set
\[
 B=\mathbb P^1\setminus\{b_0,b_1,b_2\},
 \qquad C^\circ=C\setminus\{c_1\}.
\]
The open $U_0=B\times C^\circ$ is dense in $U$.  Since $U$ is
normal and connected, $V$ and its restriction to $U_0$ are integral.

Equip $C$ with the divisorial log structure at $c_1$, and equip
$B$ with the trivial log structure.  The resulting log curve
$C^{\log}$ is proper, log smooth and log geometrically connected
over $k$, while $B$ is connected and log regular.  The primary
logarithmic product theorem gives
\[
 \pi_1(C^{\log}\times B)
   \simeq \pi_1(C^{\log})\times\pi_1(B)
   =\pi_1^t(C^\circ)\times\pi_1(B).
 \tag{*}
\]
Here we use \cite[Proposition 2.4, printed pp.~16--17]{Hoshi2006B}.
The precise identification of covers is
\cite[Proposition A.10, printed pp.~57--58]{Hoshi2006B}: for a log
regular scheme, finite Kummer \'etale covers are equivalent to
finite \'etale covers of its interior whose normalizations are tame
at every codimension-one boundary point.  Thus the sole tameness
condition in $(*)$ is along $B\times\{c_1\}$; ramification in the
$B$ direction is unrestricted.  The hypothesis at $A$ is exactly
this condition on $V|_{U_0}$, since these generic boundary valuations
agree.  Its monodromy therefore has the form
\[
 \rho(\gamma_B,\gamma_C)
     =\rho_B(\gamma_B)\rho_C(\gamma_C),
\]
where the two finite images commute, $\rho_B$ is a representation
of $\pi_1(B)$ and $\rho_C$ a representation of $\pi_1^t(C^\circ)$.
The required choices of geometric base points and paths only
conjugate these representations and do not affect vanishing of
wild inertia.

For $i=0,1$ choose a parameter $t$ at $b_i$ and an \'etale local
parameter $z$ at $c_1$.  In a blowup chart around $E_i$ we have
\[
 z=tu,\qquad E_i=\{t=0\},\qquad A=\{u=0\}.
\]
The point where $R_i$ meets $E_i$ is the point omitted from this
chart.  Choose $\lambda\in k^\times$.  The formal arc
\[
 \operatorname{Spec}k[[s]]\longrightarrow X,
 \qquad t=s,\quad u=\lambda,\quad z=\lambda s
\]
has closed point in $E_i\cap U$ and generic point in $U_0$.
The \'etale parameter on $C$ supplies the indicated lift uniquely
after fixing its closed point.  In particular the entire arc lies
in $U$.  Its pullback of $V$ is finite \'etale over the strictly
henselian ring $k[[s]]$, so its monodromy on $k((s))$ is trivial.

The $C$ component of this monodromy is tame at $s=0$: a cover
classified by $\rho_C$ is tame at $c_1$, and the parameter changes
by $z=\lambda s$.  On the wild inertia of $k((s))$, its contribution
is therefore trivial.  The product formula then forces the
$B$ component to be trivial on that same wild inertia.  Since
$t=s$, this is precisely the assertion that $\rho_B$ is tame at
$b_i$.  We have proved this for both $i=0$ and $i=1$.

At $b_2$ use the strict transform of the horizontal curve
$\mathbb P^1\times\{c_2\}$.  Its intersection with $E_2$ lies in
$U$, because $c_2\ne c_1$ and its tangent direction differs from
that of $R_2$.  An arc on this curve with parameter $t=s$ has
constant $C$ coordinate $c_2$.  Its pullback of $V$ again extends
as a finite \'etale cover of $k[[s]]$.  The $C$ component is now
trivial on the entire local fundamental group.  It follows that
$\rho_B$ is unramified at $b_2$.

Let $V_B\to B$ and $V_C\to C^\circ$ denote finite Galois covers
cut out by the kernels of $\rho_B$ and $\rho_C$, respectively.
The product cover of $U_0$ dominates $V|_{U_0}$.  At the generic
point of $R_i$, the pullback of $V_C$ is \'etale, since that point
maps into $C^\circ$.  The pullback of $V_B$ is tame, by the
conclusions at $b_i$ and tame base-change stability for discrete
valuation rings
\cite[section 15.116, Lemma 15.116.9]{StacksBIndexTame}.
Their compositum, and hence the subextension defining $V$, is
tame at $R_i$.  This proves the assertion on covers.

Finite covers of the left-hand quotient in the statement are
exactly the finite \'etale covers of $U$ tame at $A$; those of the
right-hand quotient are the covers tame at every one of the four
boundary divisors.  We have proved equality of these full
subcategories, so the natural quotient is an isomorphism.
Finally $X$ is smooth and projective, and
Proposition~\ref{D:fixed-boundary-positive} makes the right-hand
group topologically finitely generated.
\end{proof}

\begin{remark}[Scope of the horizontal reduction]
In the higher-genus version of Schr\"oer's three-point construction
\cite[sections 2.5 and 3.1--3.2, author's PDF pp.~4--5]{Schroer1999D},
the same $U$ is the proper surface with trivial Picard group minus
its three contracted points.  The proposition shows that any
finite cover which is wild at $R_0$ or $R_1$ must also be wild at
$A$.  Together with the positive-genus residue criterion, this
isolates the remaining possible covers as ones with simultaneous
horizontal wild ramification and inseparable residue extension
at a vertical exceptional divisor.  By itself this proposition does not
assert that $P_A$ has finite image in $\pi_1(U)$ or prove finite
generation of the full group.  The full group in every positive
genus is now settled by Corollary~\ref{D:schroer-all-genera},
using the separate formal-cone argument.  It uses a logarithmic product theorem with one proper log
factor; no ordinary product formula for two nonproper curves is
being assumed.
\end{remark}
